% Inserting E between P and N under concurrent writers, and the same insert
% aborting.  The companion paper's fig:listinsert adapted to the multi-writer
% case, which changes what the figure is FOR.
%
% P1's version is a four-phase lifecycle (build | install | commit | settle),
% because under exclusion the interesting thing is that a commit cannot fail and
% the reader sees one coherent list either side of it.  Here the commit CAN fail,
% so the interesting thing is the fork: one write set, one store, two outcomes.
% Drawing the two outcomes as separate figures would duplicate the shared write
% set and split the comparison across a page break; forking shows that the
% install is common and only the terminal value differs, which is the point.
%
% Two things this must make obvious that no other figure in the paper does:
%   (a) the GUARD -- a record whose old and new are the SAME value, read but not
%       written (sec:guard).  A table row with v in both columns explains
%       load-validate faster than the paragraph does.
%   (b) an abort is not a partial edit.  The FAILED branch leaves the list
%       byte-for-byte as it was found, which is what makes a losing writer's
%       retry safe.
%
% Deliberately NOT redrawn here: the slot -> record -> status anatomy, which is
% fig:tristate's job.  This figure starts from the write set and assumes it.
%
% Layout note: render the page and LOOK at it.  Every fault in this paper's other
% two figures passed a clean build.
\begin{figure}[t]
\centering
\begin{tikzpicture}[
  font=\small,
  lnode/.style   = {draw, rounded corners=2pt, minimum width=8mm, minimum height=5.5mm, fill=gray!8},
  nnode/.style   = {draw, rounded corners=2pt, minimum width=8mm, minimum height=5.5mm, fill=green!14},
  sel/.style     = {draw, very thick, minimum width=15mm, minimum height=8mm, fill=orange!18},
  trow/.style    = {anchor=west, inner sep=0pt, minimum height=6mm, font=\footnotesize},
  biarr/.style   = {latex-latex, semithick},
  arr/.style     = {-latex, thin},
  phase/.style   = {font=\footnotesize\itshape, text=black!70},
]

% ===== the write set: two writes and one guard ==============================
\node[trow, fill=blue!14] (hdr)  at (0, 3.30)
  {\makebox[15mm][c]{\textit{slot}}\makebox[12mm][c]{\textit{old}}\makebox[12mm][c]{\textit{new}}};
\node[trow, fill=blue!6]  (rec1) at (0, 2.70)
  {\makebox[15mm][c]{\code{P.next}}\makebox[12mm][c]{\code{N}}\makebox[12mm][c]{\code{E}}};
\node[trow, fill=blue!6]  (rec2) at (0, 2.10)
  {\makebox[15mm][c]{\code{N.prev}}\makebox[12mm][c]{\code{P}}\makebox[12mm][c]{\code{E}}};
\node[trow, fill=orange!10] (rec3) at (0, 1.50)
  {\makebox[15mm][c]{\code{N.next}}\makebox[12mm][c]{$v$}\makebox[12mm][c]{$v$}};
\draw[thin, blue!45] (0, 1.2) rectangle (3.9, 3.6);
\draw[thin, blue!45] (1.5, 1.2) -- (1.5, 3.6);
\draw[thin, blue!45] (2.7, 1.2) -- (2.7, 3.6);
\draw[thin, blue!45] (0, 3.0) -- (3.9, 3.0);
\draw[thin, blue!45] (0, 1.8) -- (3.9, 1.8);

% Under the table, not beside it: to the right of rec3 is where the FAILED
% branch runs, and the note printed straight over it.
\node[font=\scriptsize, text=orange!45!black] at (1.95, 0.8)
  {\textbf{guard}: read, never written};
\node[above=0.5mm of hdr, font=\footnotesize\itshape] {the write set};

% ===== the status word, forking =============================================
% Elbow arrows, not diagonals: the two branches leave the status word vertically
% and run out horizontally, which keeps their labels on a clear vertical segment
% instead of riding a slope through the outcome text.
\node[sel] (st) at (7.0, 2.4) {\code{status}};
\draw[arr] (3.9, 2.4) -- node[above, font=\scriptsize] {install} (st.west);

\draw[-latex, semithick] (7.0, 2.8) -- (7.0, 4.0) -- (9.15, 4.0);
\node[right, font=\scriptsize] at (7.05, 3.45) {\SUCCEEDED};

\draw[-latex, semithick, gray!60!black] (7.0, 2.0) -- (7.0, 0.6) -- (10.0, 0.6);
\node[right, font=\scriptsize, text=gray!45!black] at (7.05, 1.35) {\FAILED};

% ===== outcome 1: committed =================================================
% Node centres 1.7 apart against an 8mm box, so the link arrow gets a 9mm span.
% At the 1.1 spacing this started from, the gap was 3mm and two arrowheads ate
% all of it -- the boxes read as touching. Both outcome rows and all three notes
% below them share the centre 11.4, so the rows stay stacked.
\node[lnode] (cP) at ( 9.7, 4.0) {\code{P}};
\node[nnode] (cE) at (11.4, 4.0) {\code{E}};
\node[lnode] (cN) at (13.1, 4.0) {\code{N}};
\draw[biarr] (cP) -- (cE);
\draw[biarr] (cE) -- (cN);
\node[font=\scriptsize, text=green!35!black] at (11.4, 3.35)
  {\code{E} linked, both directions};

% ===== outcome 2: aborted ===================================================
\node[lnode] (aP) at (10.55, 0.6) {\code{P}};
\node[lnode] (aN) at (12.25, 0.6) {\code{N}};
\draw[biarr] (aP) -- (aN);
\node[font=\scriptsize, text=gray!45!black] at (11.4, -0.05)
  {unchanged: every record resolves to \textit{old}};
\node[font=\scriptsize, text=gray!55!black, align=center] at (11.4, -0.8)
  {a concurrent delete of \code{N} marked it,\\so the guard's \textit{old} no longer holds};

\end{tikzpicture}
\caption{One write set, one store, two outcomes---the companion paper's insert
under concurrent writers. Inserting \code{E} between \code{P} and \code{N} moves
two edges, and needs one thing more that the single-writer case did not: a
\emph{guard} on \code{N}'s forward pointer, a record whose old and new are the
same value $v$. No other record in the transaction names that slot---the insert
reads it but never writes it---so without the guard the back-edge record could
install into \code{N} before the transaction discovered, at \code{P}'s forward
pointer, that it had to abort (\cref{sec:guard}). On \SUCCEEDED\ every record resolves to its new value and
\code{E} is linked in both directions at once. On \FAILED---here because the
guard's expected value no longer holds---every record resolves to its old value
instead, so the list is left exactly as it was found and the writer may simply
re-descend and retry. An abort is never a partial edit.}
\label{fig:listinsert}
\end{figure}
